\section{引言}

过去四年，医疗人工智能经历了爆发式的快速演进。自 ChatGPT 公开发布以来，在不到三十六个月的时间里，面向医学问答的 Med-PaLM 将美国执业医师考试题集 MedQA 的准确率由 67.6\% 提升至 86.5\%[8,13]，Med-Gemini 进一步提升至 91.1\%，在修正评测题歧义后甚至达到 91.8--92.9\%[7]。同时，多模态医疗大模型也从初步接入图像，发展为能够同时处理视觉、语言、表格与基因数据的模型体系；AMIE 在 OSCE 风格诊断对话中于 30/32 个专家维度上显著优于全科医生[17]，PreA 则在两家三甲医院 2,069 名患者参与的随机对照试验中将问诊时长降低近三成，使跨科室协调度提升超过一倍[20]。

本文围绕医疗人工智能的发展图景展开，共分七节。第2节梳理 2015--2026 年医疗 AI 在训练目标、模型规模与模态广度三条主轴上的演进；第3节对比英文与中文医疗大语言模型的能力及其评估体系的变化；第4节梳理多模态医疗大模型的架构谱系；第5节论证专科基础模型在结构化电子健康记录（electronic health record, EHR）时序、影像、组学等任务上仍具有不可替代的位置；第6节讨论临床证据、评估方法与治理框架三方面尚未解决的若干问题；第7节总结全文观点。

在术语口径上，本文将医疗大语言模型（medical LLM）界定为以自然语言为主输入输出、参数规模在十亿量级以上、并在医学语料上做过领域适配（继续预训练、指令微调或强化学习）的语言模型；将多模态医疗大模型（medical LMM / VLM）界定为除文本外还能接受影像、生物信号、结构化 EHR、组学等至少一种模态作为输入的大模型；将专科基础模型界定为在单一医学子领域（如眼底、病理、CT、EHR 时序）上经自监督预训练，并能在该子领域多类下游任务上展现迁移能力的模型。医疗人工智能作为更广义的术语，涵盖上述三类以及传统任务专用监督模型。基准成绩在本文中专指静态问答或分类基准上的指标（如 MedQA 准确率），临床效用则指在真实临床工作流中以临床终点（如再入院率、误诊率、问诊时长）度量的指标。
